\section{What the Church Has Judged---and What This Study Is Doing}

\fatimathesis{Fátima is best understood neither as a codebook placed beside the Gospel nor as a picturesque residue of Portuguese folk religion. It is a historically embodied, ecclesially discerned private revelation whose durable force lies in its return to the Gospel: adore God, receive Christ, pray, repent, intercede, repair what sin has wounded, and refuse the fatalism by which violence pretends to be destiny.}

The modern reader usually approaches Fátima through a compressed sequence: three shepherd children, a dancing sun, a secret about Russia, and a papal consecration. Every item belongs to the story, but compression can falsify the whole. The six encounters unfolded between May and October 1917. Near-event interviews began while they were still occurring. The detailed account familiar today is largely Lúcia's later memoir work. The competent bishop judged the six visions worthy of belief in 1930. The first two parts of the secret were written in 1941; the third in 1944; the Holy See published the third and interpreted the whole in 2000. Pontevedra and Tuy are later experiences within a wider Fátima family, not extra pages hidden inside the 1930 decree.

\subsection{The exact ecclesial status}

On 13 October 1930 José Alves Correia da Silva, Bishop of Leiria, issued a pastoral letter after a diocesan inquiry and years of observation. Its operative clauses declare the children's visions at the Cova da Iria from 13 May through October 1917 \emph{worthy of belief} (\emph{dignas de crédito}) and officially permit the cult of Our Lady of Fátima. That is the controlling legacy judgment. It is a real positive act; it is not the 2024 category \latin{Nihil obstat}, not a modern \latin{constat de supernaturalitate} formula, and not a dogmatic definition.

The CDF's 2000 commentary supplies the assent category. An approved private revelation may be accepted with prudent human faith as credible to piety. Divine and Catholic faith is not owed to it. A Catholic may honor Our Lady, pray the Rosary, seek conversion, and live the entire sacramental life without adopting Fátima as a personal devotion. Conversely, freedom does not license misrepresentation: one should not call the event condemned, reduce the Church's decision to mere toleration, or use unresolved details to erase the actual positive judgment.

\begin{fourstable}{Reality}{Competent witness}{What may be affirmed}{Necessary limit}
Public Revelation & Christ, transmitted in Scripture and apostolic Tradition & God's definitive saving self-gift; the norm that judges every later claim & No private revelation adds a missing doctrine, sacrament, or plan of salvation.\\
The six 1917 visions & Bishop of Leiria, 1930 & Worthy of belief; cult officially permitted & Prudent human assent; no guarantee that every remembered phrase is a dictated transcript.\\
The received message and secret & Lúcia's writings under obedience; Holy See publication and commentary, 2000 & Lawful publication; a Christ-centered call to conversion; an authoritative symbolic reading & Later writing dates and visionary mediation remain material; no deterministic timetable.\\
Later Fátima cycle & Pontevedra and Tuy records; papal and pastoral reception & First Saturdays, Marian entrustment, and prayer for peoples may be received within Catholic devotion & These experiences are not silently inserted into the exact object of the 1930 clause.\\
\end{fourstable}

\Needspace{10\baselineskip}
\subsection{Four source strata, not one seamless transcript}

\begin{threestable}{Stratum}{Principal contents}{Historical force}
Near-event, 1917 & Parish interviews, Formigão and Lacerda interrogations, letters, hostile and sympathetic press, crowd testimony & Best evidence for what was said publicly and remembered near the events; incomplete, plural, and not self-authenticating.\\
Retrospective, 1922--1944 & Lúcia's first account, memoirs, correspondence, secret manuscripts & Fullest narrative and spiritual interpretation; written under obedience, but affected by elapsed time and later vocabulary.\\
Judicial and pastoral & 1922 commission, 1924 examination, 1930 letter, regulated cult & Establishes what competent authority judged and permitted; does not turn the evidence into inspired Scripture.\\
Papal and doctrinal reception & Consecrations, pilgrimages, causes of saints, CDF 2000 dossier, Benedict XVI's 2010 homily, later acts & Establishes authoritative interpretation and reception according to each act's object; cannot retroactively create a contemporaneous source.\\
\end{threestable}

The strata support one another without becoming interchangeable. A secular reporter can corroborate a crowd and a solar report without discerning its cause. A memoir can preserve a spiritually coherent whole without being contemporaneous. A bishop can judge visions credible without defining optics. A pope can receive a message providentially without canonizing every popular interpretation.

\subsection{The conclusion in advance}

Fátima's center is remarkably stable across the layers: the Rosary; adoration of the Trinity and the Eucharist; conversion from sin; prayer and willing sacrifice for sinners; peace; the Immaculate Heart as Mary's undivided \latin{fiat}; the Church passing through persecution beneath the Cross; and hope that grace can redirect freedom. Its most sensational images are ordered to the most ordinary Christian acts.

Three errors therefore fail together. \emph{Rationalistic dismissal} assumes that historical context or human mediation excludes divine action. \emph{Devotional maximalism} assumes that divine action excludes human mediation, source development, or error in expression. \emph{Apocalyptic consumerism} detaches secrets from repentance and turns prophecy into private mastery of the future. The Church's received reading is more demanding: test the testimony, respect the judgment, interpret the symbols through the Gospel, and obey the call in the present.

This study uses “reported” when the source or judgment boundary matters. It sometimes uses ordinary Catholic devotional speech---“Our Lady asked” or “the message teaches”---inside an already identified received corpus. Such usage does not promote private revelation to public Revelation or claim direct access beyond the witnesses and ecclesial acts.
